\documentclass[
	spanish,
	9pt,
	xcolor=table,
	handout,
	aspectratio=1610,
	ignorenonframetext
]{beamer}

\usepackage[spanish]{babel}
\usepackage[
	cursointroductorio,
	fullpagenumbering
]{dunestyle-beamer}
\usepackage{booktabs}
\usepackage{longtable}

\usepackage[hyperrefcolorlinks,os=win,mackeys=text]{menukeys}
\renewmenumacro{\keys}[+]{shadowedroundedkeys}

\usepackage[backend=biber,style=numeric, defernumbers=true, sorting=ynt,maxbibnames=4,maxcitenames=4]{biblatex}
\addbibresource{dune-common-references.bib}

\title{
	Introducción a línea de comandos de Linux
}
\subtitle{
	Parte I:
	\lstinline{bash},
	\lstinline{g++},
	\lstinline{cmake},
	\lstinline{gnuplot},
	\lstinline{gmsh},
	\lstinline{paraview}
}

\author{}
\institute[]
{
	\noindent
	Practique los ejemplos en Gitpod

	\href{https://gitpod.io\#/https://github.com/cpp-review-dune/dune-basics}{\includegraphics[width=4.5cm]{open-in-gitpod}}

	\vfill

	\begin{minipage}{0.12\paperwidth}
		Disponible en
	\end{minipage}%
	\hfill%
	\begin{minipage}{0.18\paperwidth}
		\href{https://github.com/cpp-review-dune}{\includegraphics[width=1.2cm]{GitHub-Mark-120px-plus}}
	\end{minipage}

	\vfill

	\begin{minipage}{0.26\paperwidth}
		¡Únete al grupo en Telegram!
	\end{minipage}%
	\hfill%
	\begin{minipage}{0.04\paperwidth}
		\href{https://t.me/joinchat/OsfYP1xnFlxjN2Ix}{\includegraphics[width=1.2cm]{telegram-logo}}
	\end{minipage}

}

\setbeamertemplate{bibliography item}{%
	\ifboolexpr{ test {\ifentrytype{book}} or test {\ifentrytype{mvbook}}
		or test {\ifentrytype{collection}} or test {\ifentrytype{mvcollection}}
		or test {\ifentrytype{reference}} or test {\ifentrytype{mvreference}} }
	{\setbeamertemplate{bibliography item}[book]}
	{\ifentrytype{online}
		{\setbeamertemplate{bibliography item}[online]}
		{\setbeamertemplate{bibliography item}[article]}}%
	\usebeamertemplate{bibliography item}}

\defbibenvironment{bibliography}
{\list{}
	{\settowidth{\labelwidth}{\usebeamertemplate{bibliography item}}%
		\setlength{\leftmargin}{\labelwidth}%
		\setlength{\labelsep}{\biblabelsep}%
		\addtolength{\leftmargin}{\labelsep}%
		\setlength{\itemsep}{\bibitemsep}%
		\setlength{\parsep}{\bibparsep}}}
{\endlist}
{\item}

\begin{document}

\frame[plain,noframenumbering]{\titlepage}

\begin{frame}
	\frametitle{Empezando con GNU/Linux}
    Nos otorga un entorno de programación con un programa base llamado shell que nos permite comunicarnos entre el kernel y las aplicaciones de usuario, por lo que controlamos el sistema operativo.
    Tiene un lenguaje de scripting que nos permite automatizar tareas que sigue el estándar POSIX.

	\begin{itemize}
		\item Abrir terminal\quad\keys{\ctrl+\Alt+t}.
        \item Obtener los comandos digitados recientemente\quad\keys{\arrowkeyup}/\keys{\arrowkeydown}.
        \item Limpiar la consola \quad\keys{\ctrl+l}.
        \item Moverse entre palabras \quad\keys{\Alt+f}/\keys{\Alt+b}.
        \item Moverse al inicio/final del comando \quad\keys{Home}/\keys{End}.
        \item Buscar comandos utilizados anteriormente\quad\keys{\ctrl+r}/\keys{\ctrl+s}.
        \item Interrumpir la ejecución de un programa\quad\keys{\ctrl+c}.
        \item Enviar a segundo plano un proceso\quad\keys{\ctrl+z}.
        \item Autocompletar el nombre de un archivo\quad\keys{\tab}.
        \item Despliega el número de posibilidades para el comando\quad\keys{\tab+\tab}.
	\end{itemize}
\end{frame}

\begin{frame}[fragile]
	\frametitle{Comandos básicos}

    Existen algunas convenciones acerca del directorio de trabajo, tiene la forma \lstinline|/foo/bar|.
    \begin{verbatim}
    [user@host somedir]$ pwd
    /home/user/somedir
    \end{verbatim}
    Cada directorio tiene por lo menos dos directorio, el directorio padre \lstinline|..| y el directorio actual \lstinline|.| como se ve a continuación.
    \begin{verbatim}
    [user@host somedir]$ ls -a
    . ..
    \end{verbatim}

	\begin{table}[ht!]
		\caption{Comandos básicos}
		\centering\footnotesize
		\begin{tabular}{cc}
			\toprule
			Comando             & Utilidad
			\tabularnewline
			\midrule
			\lstinline|pwd|  &  Muestra la ruta absoluta del directorio actual.
			\tabularnewline
            \lstinline|cd|  &  Navega al home del usuario.
			\tabularnewline
            \lstinline|ls|  &  Lista los archivos del directorio en el que se encuentra.
			\tabularnewline
            \lstinline|cp origen destino|  &  Copia un archivo origen al destino.
			\tabularnewline
            \lstinline|mkdir nombre|  &  Crea una carpeta vacía llamada nombre.
			\tabularnewline
            \lstinline|find /foo/bar -name '*.cc'|  &  Busca archivos con extensión \lstinline|.cc| en la ruta \lstinline|/foo/bar|.
			\tabularnewline
            \lstinline|chmod +x ejemplo|  &  Dar permisos de ejecución al binario ejemplo.
            \tabularnewline
            \lstinline|rm archivo|  &  Elimina el inodo del archivo.
			\tabularnewline
            \lstinline|cat archivo|  &  Imprime el contenido de un archivo a la salida estándar.
			\tabularnewline
			\bottomrule
		\end{tabular}
	\end{table}

\end{frame}

\begin{frame}
	\frametitle{Filosofía de Arch}

	\begin{description}
		\item[Simplicidad]
			Es minimalista.
			Se usa herramienta pequeñas que sigue la filosofía de UNIX, de modo que tengas una base muy pequeña que deje configurar la máquina de manera más cómoda para el usuario.
		\item[Modernidad]
			Mantener la paquetería lo más actualizada posible sin sacrificar la estabilidad. Cuenta con los compiladores más actuales de \lstinline{gcc}, \lstinline{lualatex}, \lstinline{go}, etc.
		\item[Pragmatismo]
			La gran cantidad de paquetes y scripts de compilación en los diversos repositorios de Arch Linux ofrecen software gratuito y de código abierto para quienes lo prefieren, así como paquetes de software propietario para quienes adoptan la funcionalidad por encima de la ideología.
		\item[Centrado a las usuarias y usuarios]
			La distribución está destinada a satisfacer las necesidades de quienes contribuyen a ella, en lugar de intentar atraer a tantos usuarios como sea posible.
			Está dirigido al usuario competente de GNU/Linux, o cualquier persona con una actitud de hágalo usted mismo que esté dispuesto a leer la documentación y resolver sus propios problemas.
			Se anima a todos los usuarios a participar y contribuir a la distribución.
			Informar y ayudar a corregir errores es muy valioso y los parches que mejoran los paquetes o los proyectos principales son muy apreciados: los desarrolladores de Arch son voluntarios y los contribuyentes activos a menudo se convertirán en parte de ese equipo.
		\item[Versatilidad]
	\end{description}

\end{frame}

\begin{frame}
	Estamos probando el funcionamiento de Gitpod $\sin{t}$.
	En esta ocasión hemos elegido Arch Linux como ambiente de trabajo porque tiene disponible una gran variedad de módulos de DUNE.
	Es recomendable habilitar el repositorio \href{https://wiki.archlinux.org/title/Unofficial\_user\_repositories\#arch4edu}{\texttt{arch4edu}}\footnote{Administrado por Jingbei Li de la Universidad de Tsinghua.}.
\end{frame}

\begin{frame}
	\frametitle{Usando el emulador de terminal}
	% Una consola es un display del sistema
	% En todo lo que sigue, usaremos el shell \lstinline{bash}.

	% GNU General Public License v2.0
	% GNU Lesser General Public License v2.1
	% BSD 3-Clause "New" or "Revised" License

	% El terminal te da acceso al shell, que nos da acceso al sistema operativo subyacente.
	% El shell, sabe interactuar entre el usuario y el sistema operativo.

	% CLI Command Line Interface
\end{frame}

\begin{frame}
	% pwd
	% El cursor parpadeante al final de la línea es donde se ingresan los comandos basados ​​en texto. Puede comenzar a experimentar con un comando simple, Imprimir directorio de trabajo (pwd), que muestra la carpeta actual en la que se encuentra en la pantalla. Escriba pwd y presione Entrar.

	% Todos los comandos que ingresa funcionan de la misma manera: ingresa el comando, incluye cualquier parámetro para extender el uso del comando y presiona Enter para ejecutar la línea de comando que ingresó. Escribe en la Terminal: uname -a y presiona Enter. Esto muestra información del sistema con respecto a Mint
\end{frame}

\begin{frame}[fragile]
	\begin{lstlisting}
    gitpod ~/dune-basics $ cd
  \end{lstlisting}
\end{frame}

%\lstinputpath{../../../sandbox}

\begin{frame}[fragile]
	\frametitle{Classes}

	\lstinputlisting[
		caption={Programa \texttt{hello-linux.cc}.},
		label=hello-linux.cc,
	]{../../../sandbox/hello-linux.cc}


\end{frame}

\begin{frame}[fragile]
	\frametitle{Classes}

	\lstinputlisting[
		caption={Programa \texttt{dune-basics.cc}.},
		label=dune-basics.cc,
	]{../../../src/dune-basics.cc}

\end{frame}

\begin{frame}[fragile]
	\frametitle{Classes}

	\lstinputlisting[
		caption={Programa \texttt{dune-basics.cc}.},
		label=dune-basics.cc,
	]{../../../sandbox/dune-math-constants.cc}

\end{frame}
%https://archlinux.org/logos/archlinux-icon-crystal-256.svg

\begin{frame}
	\frametitle{El comando \lstinline{duneproject}}
	% https://gitlab.dune-project.org/core/dune-common/-/raw/master/bin/duneproject
	Es un asistente en el lenguaje \lstinline{bash} que se encuentra en \lstinline{/usr/bin/duneproject}\ldots

\end{frame}

\begin{frame}\transblindsvertical
	\frametitle{Referencias}
	%------------------------------------------------------------ 1
	\only<1>{
		\begin{itemize}
			\item Libros
			      \nocite{*}
			      \printbibliography[heading=none,keyword=book]
		\end{itemize}
	}
	%------------------------------------------------------------ 2
	\only<2>{
		\begin{itemize}
			\item Artículos
			      \printbibliography[heading=none,keyword=paper]
		\end{itemize}
	}
	%------------------------------------------------------------ 3
	\only<3>{
		\begin{itemize}
			\item Sitios web
			      \printbibliography[heading=none,keyword=online]
		\end{itemize}
	}
\end{frame}

\end{document}
%https://github.com/easybuilders/easybuild/wiki/EasyBuild-Tech-Talks-I%3A-Open-MPI
% https://www.youtube.com/playlist?list=PLagFkXs2BczZO4B_I6YhuS4xOxlSPgZuT
% https://www.youtube.com/playlist?list=PLagFkXs2BczZGsvnN7UAqd5LWfgOrat0v
% https://www.youtube.com/playlist?list=PLZRRlbOTxTmAMASXs7mAFnKkFZ3ohefTF
% https://www.youtube.com/playlist?list=PLZRRlbOTxTmBERH5Ov1kWhPlmOQAP04zi
% https://www.youtube.com/playlist?list=PLZRRlbOTxTmCaBYeLHQyqMFyI6caKsVDd
% https://www.youtube.com/playlist?list=PLZRRlbOTxTmB_gV8rLQ6N9KDX5po5eTay
% https://www.youtube.com/playlist?list=PLZRRlbOTxTmCyLkmjWg0e8cNk6rONT-mk